%-----------------------------------------------------------------------------------------------------------------------------------------------%
%	The MIT License (MIT)
%
%	Copyright (c) 2019 Jan Küster
%
%	Permission is hereby granted, free of charge, to any person obtaining a copy
%	of this software and associated documentation files (the "Software"), to deal
%	in the Software without restriction, including without limitation the rights
%	to use, copy, modify, merge, publish, distribute, sublicense, and/or sell
%	copies of the Software, and to permit persons to whom the Software is
%	furnished to do so, subject to the following conditions:
%	
%	THE SOFTWARE IS PROVIDED "AS IS", WITHOUT WARRANTY OF ANY KIND, EXPRESS OR
%	IMPLIED, INCLUDING BUT NOT LIMITED TO THE WARRANTIES OF MERCHANTABILITY,
%	FITNESS FOR A PARTICULAR PURPOSE AND NONINFRINGEMENT. IN NO EVENT SHALL THE
%	AUTHORS OR COPYRIGHT HOLDERS BE LIABLE FOR ANY CLAIM, DAMAGES OR OTHER
%	LIABILITY, WHETHER IN AN ACTION OF CONTRACT, TORT OR OTHERWISE, ARISING FROM,
%	OUT OF OR IN CONNECTION WITH THE SOFTWARE OR THE USE OR OTHER DEALINGS IN
%	THE SOFTWARE.
%	
%
%-----------------------------------------------------------------------------------------------------------------------------------------------%

%============================================================================%
%
%	DOCUMENT DEFINITION
%
%============================================================================%

%we use article class because we want to fully customize the page and don't use a cv template
\documentclass[8pt,A4]{article}	


%----------------------------------------------------------------------------------------
%	ENCODING
%----------------------------------------------------------------------------------------

% we use utf8 since we want to build from any machine
\usepackage[utf8]{inputenc}		

%----------------------------------------------------------------------------------------
%	LOGIC
%----------------------------------------------------------------------------------------

% provides \isempty test
\usepackage{xstring, xifthen}

%----------------------------------------------------------------------------------------
%	FONT BASICS
%----------------------------------------------------------------------------------------

% some tex-live fonts - choose your own

%\usepackage[defaultsans]{droidsans}
%\usepackage[default]{comfortaa}
%\usepackage{cmbright}
\usepackage[default]{raleway}
%\usepackage{fetamont}
%\usepackage[default]{gillius}
%\usepackage[light,math]{iwona}
%\usepackage[thin]{roboto} 

% set font default
\renewcommand*\familydefault{\sfdefault} 	
\usepackage[T1]{fontenc}

% more font size definitions
\usepackage{moresize}

% redefine itemss
\usepackage{enumitem}
\setlist[itemize]{noitemsep, topsep=2pt}

%----------------------------------------------------------------------------------------
%	FONT AWESOME ICONS
%---------------------------------------------------------------------------------------- 

% include the fontawesome icon set
\usepackage{fontawesome}

% use to vertically center content
% credits to: http://tex.stackexchange.com/questions/7219/how-to-vertically-center-two-images-next-to-each-other
\newcommand{\vcenteredinclude}[1]{\begingroup
\setbox0=\hbox{\includegraphics{#1}}%
\parbox{\wd0}{\box0}\endgroup}

% use to vertically center content
% credits to: http://tex.stackexchange.com/questions/7219/how-to-vertically-center-two-images-next-to-each-other
\newcommand*{\vcenteredhbox}[1]{\begingroup
\setbox0=\hbox{#1}\parbox{\wd0}{\box0}\endgroup}

% icon shortcut
\newcommand{\icon}[3] { 							
	\makebox(#2, #2){\textcolor{maincol}{\csname fa#1\endcsname}}
}	

% icon with text shortcut
\newcommand{\icontext}[4]{ 						
	\vcenteredhbox{\icon{#1}{#2}{#3}}  \hspace{2pt}  \parbox{0.9\mpwidth}{\textcolor{#4}{#3}}
}

% icon with website url
\newcommand{\iconhref}[5]{ 						
    \vcenteredhbox{\icon{#1}{#2}{#5}}  \hspace{2pt} \href{#4}{\textcolor{#5}{#3}}
}

% icon with email link
\newcommand{\iconemail}[5]{ 						
    \vcenteredhbox{\icon{#1}{#2}{#5}}  \hspace{2pt} \href{mailto:#4}{\textcolor{#5}{#3}}
}

%----------------------------------------------------------------------------------------
%	PAGE LAYOUT  DEFINITIONS
%----------------------------------------------------------------------------------------

% page outer frames (debug-only)
% \usepackage{showframe}		

% we use paracol to display breakable two columns
\usepackage{paracol}

% define page styles using geometry
\usepackage[a4paper]{geometry}

% remove all possible margins
\geometry{top=1cm, bottom=1cm, left=1cm, right=1cm}

\usepackage{fancyhdr}
\pagestyle{empty}

% space between header and content
% \setlength{\headheight}{0pt}

% indentation is zero
\setlength{\parindent}{0mm}

%----------------------------------------------------------------------------------------
%	TABLE /ARRAY DEFINITIONS
%---------------------------------------------------------------------------------------- 

% extended aligning of tabular cells
\usepackage{array}

% custom column right-align with fixed width
% use like p{size} but via x{size}
\newcolumntype{x}[1]{%
>{\raggedleft\hspace{0pt}}p{#1}}%


%----------------------------------------------------------------------------------------
%	GRAPHICS DEFINITIONS
%---------------------------------------------------------------------------------------- 

%for header image
\usepackage{graphicx}

% use this for floating figures
% \usepackage{wrapfig}
% \usepackage{float}
% \floatstyle{boxed} 
% \restylefloat{figure}

%for drawing graphics		
\usepackage{tikz}				
\usetikzlibrary{shapes, backgrounds,mindmap, trees}

%----------------------------------------------------------------------------------------
%	Color DEFINITIONS
%---------------------------------------------------------------------------------------- 
\usepackage{transparent}
\usepackage{color}

% primary color
% \definecolor{maincol}{RGB}{ 18, 78, 137 }
\definecolor{maincol}{RGB}{ 18, 135, 137}


% accent color, secondary
 \definecolor{accentcol}{RGB}{ 250, 150, 10 }

% dark color
\definecolor{darkcol}{RGB}{25, 60, 62}

% light color
\definecolor{lightcol}{RGB}{245,245,245}


% Package for links, must be the last package used
\usepackage[hidelinks]{hyperref}

% returns minipage width minus two times \fboxsep
% to keep padding included in width calculations
% can also be used for other boxes / environments
\newcommand{\mpwidth}{\linewidth-\fboxsep-\fboxsep}
	


%============================================================================%
%
%	CV COMMANDS
%
%============================================================================%

%----------------------------------------------------------------------------------------
%	 CV LIST
%----------------------------------------------------------------------------------------

% renders a standard latex list but abstracts away the environment definition (begin/end)
\newcommand{\cvlist}[1] {
	\begin{itemize}{\small #1}\end{itemize}
}

%----------------------------------------------------------------------------------------
%	 CV TEXT
%----------------------------------------------------------------------------------------

% base class to wrap any text based stuff here. Renders like a paragraph.
% Allows complex commands to be passed, too.
% param 1: *any
\newcommand{\cvtext}[1] {
	\begin{tabular*}{1\mpwidth}{p{0.98\mpwidth}}
		\parbox{1\mpwidth}{\small #1}
	\end{tabular*}
}

%----------------------------------------------------------------------------------------
%	CV SECTION
%----------------------------------------------------------------------------------------

% Renders a a CV section headline with a nice underline in main color.
% param 1: section title
\newcommand{\cvsection}[1] {
	\vspace{14pt}
	\cvtext{
		\textbf{\LARGE{\textcolor{darkcol}{\uppercase{#1}}}}\\[-4pt]
		\textcolor{maincol}{ \rule{0.1\textwidth}{2pt} } \\
	}
}

%----------------------------------------------------------------------------------------
%	META SKILL
%----------------------------------------------------------------------------------------

% Renders a progress-bar to indicate a certain skill in percent.
% param 1: name of the skill / tech / etc.
% param 2: level (for example in years)
% param 3: percent, values range from 0 to 1
\newcommand{\cvskill}[3] {
	\begin{tabular*}{1\mpwidth}{p{0.72\mpwidth}  r}
 		\textcolor{black}{\textbf{#1}} & \textcolor{maincol}{#2}\\
	\end{tabular*}%
	
	\hspace{4pt}
	\begin{tikzpicture}[scale=1,rounded corners=2pt,very thin]
		\fill [lightcol] (0,0) rectangle (1\mpwidth, 0.15);
		\fill [maincol] (0,0) rectangle (#3\mpwidth, 0.15);
  	\end{tikzpicture}%
}


%----------------------------------------------------------------------------------------
%	 CV EVENT
%----------------------------------------------------------------------------------------

% Renders a table and a paragraph (cvtext) wrapped in a parbox (to ensure minimum content
% is glued together when a pagebreak appears).
% Additional Information can be passed in text or list form (or other environments).
% the work you did
% param 1: time-frame i.e. Sep 14 - Jan 15 etc.
% param 2:	 event name (job position etc.)
% param 3: Customer, Employer, Industry
% param 4: Short description
% param 5: work done (optional)
% param 6: technologies include (optional)
% param 7: achievements (optional)
\newcommand{\cvevent}[7] {
	
	% we wrap this part in a parbox, so title and description are not separated on a pagebreak
	% if you need more control on page breaks, remove the parbox
	\parbox{\mpwidth}{
		\begin{tabular*}{1\mpwidth}{p{0.72\mpwidth}  r}
	 		\textcolor{black}{\textbf{#2}} & \colorbox{maincol}{\makebox[0.25\mpwidth]{\textcolor{white}{#1}}} \\
			\textcolor{maincol}{\textbf{#3}} & \\
		\end{tabular*}\\[8pt]
	
		\ifthenelse{\isempty{#4}}{}{
			\cvtext{#4}\\
		}
	}

	\ifthenelse{\isempty{#5}}{}{
		\vspace{9pt}
		{#5}
	}

	\ifthenelse{\isempty{#6}}{}{
		\vspace{9pt}
		\cvtext{\textbf{Technologies utilisées:}}\\
		{#6}
	}

	% \ifthenelse{\isempty{#7}}{}{
	% 	\vspace{9pt}
	% 	\cvtext{\textbf{Achievements include:}}\\
	% 	{#7}
	% }
	\vspace{14pt}
}

%----------------------------------------------------------------------------------------
%	 CV META EVENT
%----------------------------------------------------------------------------------------

% Renders a CV event on the sidebar
% param 1: title
% param 2: subtitle (optional)
% param 3: customer, employer, etc,. (optional)
% param 4: info text (optional)
\newcommand{\cvmetaevent}[4] {
	\textcolor{maincol} {\cvtext{\textbf{\begin{flushleft}#1\end{flushleft}}}}

	\ifthenelse{\isempty{#2}}{}{
	\textcolor{darkcol} {\cvtext{\textbf{#2}} }
	}

	\ifthenelse{\isempty{#3}}{}{
		\cvtext{{ \textcolor{darkcol} {#3} }}\\
	}

	\cvtext{#4}\\[14pt]
}

%---------------------------------------------------------------------------------------
%	QR CODE
%----------------------------------------------------------------------------------------

% Renders a qrcode image (centered, relative to the parentwidth)
% param 1: percent width, from 0 to 1
\newcommand{\cvqrcode}[1] {
	\begin{center}
		
\includegraphics[width={#1}\mpwidth]{latex/qrcode}
	\end{center}
}


%============================================================================%
%
%
%
%	DOCUMENT CONTENT
%
%
%
%============================================================================%
\begin{document}
\columnratio{0.28}
\setlength{\columnsep}{2.2em}
\setlength{\columnseprule}{4pt}
\colseprulecolor{lightcol}
\begin{paracol}{2}
	\begin{leftcolumn}
		%---------------------------------------------------------------------------------------
		%	META IMAGE
		%----------------------------------------------------------------------------------------
		\includegraphics[width=\linewidth]{latex/bruno_avatar.png}	%trimming relative to image size

		%---------------------------------------------------------------------------------------
		%	META SKILLS
		%----------------------------------------------------------------------------------------
		\cvsection{Compétences}

		\cvskill{Linux} {20+ yrs} {1} \\[-2pt]
		\cvskill{Python} {15+ yrs} {0.9} \\[-2pt]
		\cvskill{AWS} {10+ yrs} {0.9} \\[-2pt]
		\cvskill{Terraform} {10+ yrs} {0.8} \\[-2pt]
		\cvskill{Kubernetes} {10+ yrs} {0.8} \\[-2pt]
		\cvskill{Helm} {10+ yrs} {0.8} \\[-2pt]
		\cvskill{Go} {3+ yrs} {0.6} \\[-2pt]
		\cvskill{TypeScript} {2+ yrs} {0.4} \\[-2pt]
		\cvskill{NixOS} {4+ yrs} {0.8} \\[-2pt]

		\vfill\null
		\cvsection{CONTACT}

		\icontext{MapMarker}{12}{Montpellier\\France}{black}\\[6pt]
		\icontext{MobilePhone}{12}{+33 6 14 62 06 95}{black}\\[6pt]
		\iconemail{Envelope}{12}{brunoadele@gmail.com}{brunoadele@gmail.com}{black}\\[6pt]
		\icontext{Github}{12}{github.com/badele}{black}\\[6pt]

		\vfill\null
		\cvqrcode{0.7}

		%---------------------------------------------------------------------------------------
		%	EDUCATION
		%----------------------------------------------------------------------------------------
		\newpage
		\cvsection{LANGUES}

		\vspace{-24pt}

		\cvmetaevent
		{}
		{Francais}
		{Natif}

		\vspace{-64pt}

		\cvmetaevent
		{}
		{Anglais}
		{Lu / Ecrit}

		\cvsection{HOBBIES}

		\vspace{-24pt}

		\begin{itemize}
			\item Informatique/Domotique
			\item Image de synthèse
			\item Photo
			\item Vélo/Badminton/Rando
		\end{itemize}

		\cvsection{PROJETS}

		\icontext{Star}{12}{\textbf{988} Total GitHub stars}{black}\\[6pt]

\icontext{Star}{12}{\textbf{479}  \href{https://github.com/badele/nix-homelab}{badele/nix-homelab}}{black}\\[6pt]
\icontext{Star}{12}{\textbf{47}  \href{https://github.com/badele/vide}{badele/vide}}{black}\\[6pt]
\icontext{Star}{12}{\textbf{27}  \href{https://github.com/badele/network-lab}{badele/network-lab}}{black}\\[6pt]
\icontext{Star}{12}{\textbf{12}  \href{https://github.com/badele/slide-generator}{badele/slide-generator}}{black}\\[6pt]
\icontext{Star}{12}{\textbf{2}  \href{https://github.com/badele/splitans}{badele/splitans}}{black}\\[6pt]


		\vfill\null
		\cvqrcode{0.7}

	\end{leftcolumn}
	\begin{rightcolumn}
		%---------------------------------------------------------------------------------------
		%	TITLE  HEADER
		%----------------------------------------------------------------------------------------
		% \fcolorbox{white}{darkcol}{\begin{minipage}[c][3.5cm][c]{1\mpwidth}
		% 		\begin {center}
		% 		\HUGE{ \textbf{ \textcolor{white}{ \uppercase{ Bruno Adelé } } } } \\[-24pt]
		% 		\textcolor{white}{ \rule{0.1\textwidth}{1.25pt} } \\[4pt]
		% 		\large{ \textcolor{white} {Senior DevOps / SRE} }
		% 		\end {center}
		% 	\end{minipage}} \\[14pt]
		% \vspace{-18pt}
		\fcolorbox{white}{darkcol}{%
			\begin{minipage}[c][3.5cm][c]{1\mpwidth}
				\begin{center}
					\HUGE{\textbf{\textcolor{white}{\uppercase{Bruno Adelé}}}}%
					\vspace{-24pt}

					\textcolor{white}{\rule{0.1\textwidth}{1.25pt}}

					\large{\textcolor{white}{Senior DevOps / SRE}}
				\end{center}
			\end{minipage}%
		}

		\vspace{14pt}
		%---------------------------------------------------------------------------------------
		%	PROFILE
		%----------------------------------------------------------------------------------------
		\vfill\null
		\cvsection{PROFIL}

		\cvtext{Passionné d’informatique depuis mon plus jeune âge, je possède \textbf{28 ans d’expérience} sous \textbf{Linux}, dont \textbf{10 ans} sur le cloud (\textbf{8 ans sur AWS / 2 ans sur GCP}). J’ai une forte inclination pour les infrastructures \textbf{immuables} et une volonté de \textbf{maximiser les performances} tout en \textbf{minimisant les ressources et les coûts}. J’aime construire des architectures et aider les développeurs à optimiser leurs applications et leurs déploiements.\\

			Je suis égallement attentif à la q\textbf{ualité de service délivrée} aux \textbf{utilisateurs finaux}, avec une approche orientée fiabilité, disponibilité et amélioration continue.\\

		}

		%---------------------------------------------------------------------------------------
		%	WORK EXPERIENCE
		%----------------------------------------------------------------------------------------
		\vfill\null
		\cvsection{Experiences}

		\cvevent
		{2024 - 2026}
		{Senior Lead Devops / SRE}
		{Araiko, Lyon}
		{DevOps / SRE en charge de la gestion et de la \textbf{fiabilité} de la plateforme : pilotage des infrastructures AWS (EKS), des pipelines CI/CD et de la supervision. \textbf{Référent technique} pour les développeurs et encadrement d’un ingénieur DevOps.}
		{\cvlist{
				\item Gestion et automatisation d’infrastructures \textbf{AWS} multi-environnements via \textbf{Terraform}
				\item Développement de \textbf{pipelines CI/CD} avec GitHub Actions pour fiabiliser et simplifier les déploiements des équipes de développement
				\item Mise en place de la supervision (\textbf{Prometheus}, \textbf{Grafana})
				\item pilotage de l’\textbf{optimisation des coûts} afin d’améliorer la performance et l’efficience de la plateforme
			}}
		{\cvlist {
				\item \textbf{AWS}, \textbf{EKS}, \textbf{Kubernetes}, \textbf{Docker} pour la gestion de l'infrastructure
				\item \textbf{Terraform}, \textbf{Helm}, \textbf{Python}, \textbf{Typescript}, \textbf{Nix}, \textbf{Just}  pour l'automatisation
				\item \textbf{Prometheus}, \textbf{Grafana} pour la supervision de la plateforme
			}}

		\vfill\null
		\cvevent
		{2021 - 2024}
		{Senior Devops / SRE}
		{SeqOne Genomics, Montpellier}
		{DevOps / SRE garant de la stabilité et des performances de la plateforme : amélioration continue de la fiabilité via la supervision. Conception d’architectures cloud \textbf{immuables} optimisées en coûts et ressources, et déploiement des applications sur \textbf{Kubernetes} avec \textbf{Helm}.}
		{\cvlist{
				\item Gestion et amélioration continue de la fiabilité de la plateforme
				\item Conception d’architectures cloud \textbf{immuables optimisées en coûts et ressources}
				\item Déploiement et industrialisation des applications sur \textbf{Kubernetes} avec \textbf{Helm}
				\item Mise en place et exploitation de la supervision (\textbf{Prometheus}, \textbf{Grafana})
			}}
		{\cvlist {
				\item \textbf{AWS}, \textbf{EKS}, \textbf{Kubernetes}, \textbf{Docker} pour la gestion de l'infrastructure
				\item \textbf{Terraform}, \textbf{Helm}, \textbf{Python} pour l'automatisation
				\item \textbf{Prometheus}, \textbf{Grafana} pour la supervision de la plateforme
			}}

		\vfill\null
		\cvevent
		{2018 - 2021}
		{Senior Devops / SRE}
		{Adagio, Montpellier}
		{DevOps / SRE en charge de la mise en place et de la fiabilité des infrastructures cloud : conception et automatisation sur \textbf{GCP} via Infrastructure as Code, et industrialisation des déploiements applicatifs.}
		{\cvlist{
				\item Conception et automatisation d’infrastructures \textbf{GCP} via \textbf{Terraform}
				\item Mise en place de pipelines CI/CD avec GitHub Actions et GitLab CI
				\item Déploiement et gestion d’applications conteneurisées (\textbf{Docker})
				\item Mise en place de la supervision et du monitoring (\textbf{Datadog})
			}}
		{\cvlist {
				\item \textbf{GCP}, \textbf{Docker} pour l'infrastructure et les déploiements
				\item \textbf{Terraform}, \textbf{Puppet} pour l'automatisation et la gestion de configuration
				\item \textbf{GitHub Actions}, \textbf{GitLab CI }pour les pipelines CI/CD
				\item Datadog pour la supervision
			}}

		\vfill\null
		\cvevent
		{2015 - 2018}
		{Senior Devops / SRE}
		{Hawk, Montpellier}
		{DevOps / SRE en charge de la conception et de la fiabilité d’infrastructures multi-cloud (\textbf{AWS, GCP, OVH}) à forte volumétrie : sécurisation des interconnexions, mise en place de la supervision et exploitation d’une plateforme supportant jusqu’à \textbf{300\,000 requêtes} par seconde.}
		{\cvlist{
				\item Conception et automatisation d’infrastructures multi-cloud (\textbf{AWS, GCP})
				\item Mise en place d’interconnexions sécurisées via \textbf{IPSec}
				\item Déploiement et exploitation de la supervision (\textbf{Consul}, \textbf{Telegraf}, \textbf{InfluxDB}, \textbf{Grafana})
				\item Gestion d’une plateforme à forte charge (\textbf{300\,000 req/s})
			}}
		{\cvlist {
				\item \textbf{AWS}, \textbf{GCP} pour l'infrastructure cloud
				\item \textbf{Terraform}, \textbf{Salt} pour l'automatisation
				\item \textbf{IPSec} pour la sécurisation des interconnexions
				\item \textbf{Consul}, \textbf{Telegraf}, \textbf{InfluxDB}, \textbf{Grafana} pour la supervision
			}}

		\vfill\null
		\cvevent
		{1998 - 2015}
		{Architecte IT \& Systèmes d'Information}
		{CompuGroup Medical AG, Montpellier}
		{Architecte IT en charge de la conception, du déploiement et de l’évolution des systèmes d’information : définition des architectures, \textbf{pilotage des budgets} et \textbf{supervision des infrastructures}.}
		{\cvlist{
				\item Conception et déploiement d’architectures virtualisées
				\item Pilotage budgétaire des infrastructures IT
				\item Supervision et maintien en conditions opérationnelles des systèmes
			}}
		{\cvlist {
				\item \textbf{Virtualisation}, infrastructures systèmes
				\item Outils de supervision
			}}


		% hotfixes to create fake-space to ensure the whole height is used
		\mbox{}
		\vfill
		\mbox{}
		\vfill
		\mbox{}
		\vfill
		\mbox{}
	\end{rightcolumn}
\end{paracol}
\end{document}
